\chapter{怪词速查（零基础版）}
\label{ch:N-jargon}

按拼音/英文混排，遇到本册任何怪词先查这里。

\begin{description}[style=nextline,leftmargin=2.2cm]
  \item[avg / mean] 希望块的平均厚度。
  \item[beta1] Chain：骨架有没有环。
  \item[calib] 试切后写下的调弦结果。
  \item[CDC] 按内容找刀口的分块方法。
  \item[chunk] 一块连续字节。
  \item[Content ID] 块上的指纹贴纸（常用 SHA256）。
  \item[dedup] 相同贴纸只存一份。
  \item[Embed / Y] Native 用邻近字节炼出的星号。
  \item[feed] 流式喂入的一大盒数据（常 512KB）。
  \item[flag] 仓库房门上的类型牌子。
  \item[Flip] 星图朝向翻转计数。
  \item[Gear] 一种滚动指纹搅拌机。
  \item[hard cut] 走到 max 仍无信号则强制切。
  \item[hash] 把数据打成摘要数字/贴纸。
  \item[Hom] 第三代：铃 + 珠串变化。
  \item[landmark] 采样星点。
  \item[lazy] 没叮就不做贵操作。
  \item[LFSR] Chain 的密码锁状态机。
  \item[mask] 只检查指纹某些位。
  \item[min / max] 最短/最长块约束。
  \item[Native] 第五·一代：星图事件即刀口。
  \item[parity] 整文件切点等于分盒拼接切点。
  \item[PH / H0] 看连通团数量变化。
  \item[primary] 第一关（铃或对齐）。
  \item[quant] 把字节颜色弄粗糙。
  \item[resume] 分盒之间带走的进度小条。
  \item[reuse] 内容平移后刀口仍能对上的比例。
  \item[seed] 菜谱编号，决定表/锁/炼金。
  \item[SIMD] 一次算多格的加速外壳。
  \item[stride] 网格间距/对齐稀密度。
  \item[TopoPH] 第五代混合：铃 + 滤网。
  \item[window] 眼前固定宽度的字节段。
\end{description}

\vspace{0.8em}
\noindent\textit{查完词若仍懵：回到第二部分动画片，用剧本句子套词。}
